\subsection*{MT-M12. An exact simple-zero gate for the canonical arithmetic quotient}
Fix a primitive unitary Hecke character $\eta$ over a number field
$K$, and assume GRH for its $L$-function. Assume additionally that
all its nontrivial zeros are simple. Let
\[
 \Xi_\eta(s)=P_\eta(s)\Lambda(s,\eta)
\]
be the entire completed function with its possible poles removed.
For a non-pure-norm character take $P_\eta=1$. If
$L(s,\eta)=\zeta_K(s+i\tau)$, take
$P_\eta(s)=(s+i\tau)(s-1+i\tau)$.
The zeros of $\Xi_\eta$ are precisely the nontrivial zeros, denoted
$Z$. Put
\[
 \psi_\eta(x)=\frac1{2\pi}\int_{\mathbb R}\Xi_\eta(it)e^{-itx}\,dt,
 \qquad Q_\eta=\mathcal E/\overline{\psi_\eta*\mathcal E}.
\]
Then $\psi_\eta\in\mathcal E$, with
$\mathcal M\psi_\eta=\Xi_\eta$, and the following are equivalent:
\begin{enumerate}
\item There are $c,C>0$ such that
$|L'(\rho,\eta)|\ge c(1+|\operatorname{Im}\rho|)^{-C}$ for
every $\rho\in Z$.
\item $Q_\eta$ has an absolute Schauder basis with coefficient
functionals $[f]\mapsto\mathcal Mf(\rho)$, $\rho\in Z$.
\item $Q_\eta$ has a Schauder basis with these continuous
coefficient functionals.
\end{enumerate}
Under these conditions the closed convolution ideal is exactly
\[
 K_Z=\{f\in\mathcal E:\mathcal Mf(\rho)=0
                              \text{ for every }\rho\in Z\}. \tag{MT-M12a}
\]
Moreover $Q_\eta$ is topologically the rapidly decreasing sequence
space
\[
 s(Z)=\{(a_\rho):\sup_{\rho\in Z}
       (1+|\operatorname{Im}\rho|)^k|a_\rho|<\infty
                                  \text{ for every }k\ge0\}. \tag{MT-M12b}
\]
This equivalence uses $A_D$ for Dirichlet functions, or $M_H$
conditionally for general unitary Hecke functions. Simplicity is
an extra assumption throughout; condition (1) is the remaining
gate, not a consequence of those GRH inputs proved here.

\begin{proof}
\emph{Analytic preliminaries.}
Write $n=[K:\mathbb Q]$ and $\kappa=\pi n/4$.
The fixed-function polynomial strip bounds of NT-Z10 and
Stirling's formula for the $n$ factors $\Gamma_{\mathbb R}$
give, for every fixed $A>0$, constants $C_A,B_A$ such that
\[
 |\Xi_\eta(\sigma+it)|
       \le C_A(1+|t|)^{B_A}e^{-\kappa|t|}
                            \qquad(|\sigma|\le A).     \tag{5}
\]
On bounded $t$ the entire function is bounded on each compact
strip rectangle, so enlarging $C_A$ covers possible gamma-pole
cancellations there. Derivative bounds follow by Cauchy's formula
on a slightly larger strip. Thus $\Xi_\eta(it)$ is Schwartz.
The contour-shifting and Fourier-inversion argument of MT-M09,
using (5) in place of the Gaussian bound, shows that the displayed
inverse transform belongs to $\mathcal E$ and has Mellin transform
$\Xi_\eta$. More explicitly, shift to $\sigma=N+1$ on $x\ge0$
and to $\sigma=-N-1$ on $x\le0$; every integral with the extra
factor $(-\sigma-it)^j$ is absolutely convergent by (5).
This proves all $p_{N,j}$ bounds. The horizontal contour edges
tend to zero by the same estimate, and Fourier uniqueness then
proves the transform identity.

Assume condition (1). At a zero $\rho=1/2+iT$, differentiation
of the completed function gives
\[
 \Xi_\eta'(\rho)
  =P_\eta(\rho)q_\eta^{\rho/2}
       \prod_j\Gamma_{\mathbb R}(\rho+\mu_j)L'(\rho,\eta),
\]
where $q_\eta=|D_K|\operatorname N\mathfrak f_\eta$
is the finite conductor.
Stirling's formula and (1) imply, after enlarging constants to
cover the finitely many low zeros,
\[
        |\Xi_\eta'(\rho)|
           \ge c_1(1+|T|)^{-C_1}e^{-\kappa|T|}.         \tag{6}
\]

\emph{Polynomially bounded cardinal functions.}
Fix $a>0$. For each zero $\rho$ define the entire function
\[
 H_\rho(s)=
     \frac{\Xi_\eta(s)}
          {\Xi_\eta'(\rho)(s-\rho)}\,e^{a(s-\rho)^2}.    \tag{7}
\]
It takes the value one at $\rho$ and zero at every other zero.
For each fixed strip $|\operatorname{Re}s|\le A$, (5)--(6)
imply a bound
\[
 |H_\rho(\sigma+it)|
       \le C_A'(1+|T|)^{B_A'}
                              e^{-a(t-T)^2/2}.         \tag{8}
\]
Here are the details near the apparent denominator. If
$|s-\rho|\ge1$, use (5) and (6), together with
$\big||T|-|t|\big|\le|T-t|$ and
$1+|t|\le(1+|T|)(1+|t-T|)$.
The remaining fixed polynomial and
$e^{\kappa|t-T|}$ are absorbed by half the Gaussian.
If $|s-\rho|<1$, use
\[
 \frac{\Xi_\eta(s)}{s-\rho}
       =\int_0^1\Xi_\eta'(\rho+u(s-\rho))\,du
\]
and the Cauchy derivative bound from (5) on a slightly larger
strip. Since $|t-T|<1$, the same estimate (8) follows by
adjusting the constants.

Invert $H_\rho$ by Fourier transform and call the inverse
$h_\rho$. The same contour argument gives $h_\rho\in\mathcal E$,
$\mathcal Mh_\rho=H_\rho$, and, for each $N,j$,
\[
             p_{N,j}(h_\rho)
                    \le C_{N,j}(1+|T|)^{B_{N,j}}.      \tag{9}
\]
Indeed, on the shifted contour used for $x\ge0$ or $x\le0$,
integrating $(1+|t|)^j$ against the Gaussian in (8) costs
only another fixed power of $1+|T|$.

For any $f\in\mathcal E$, its values $\mathcal Mf(1/2+iT)$
decrease faster than every power of $|T|$. This follows by
integrating by parts arbitrarily many times in
$\int e^{x/2}f(x)e^{iTx}\,dx$; each resulting $L^1$ norm
is bounded by finitely many defining seminorms of $f$.
The local count NT-Z10 implies
$\sum_{\rho\in Z}(1+|\operatorname{Im}\rho|)^{-2}<\infty$.
Consequently
\[
              I(f)=\sum_{\rho\in Z}\mathcal Mf(\rho)h_\rho
                                                               \tag{10}
\]
converges absolutely in every seminorm of $\mathcal E$,
defines a continuous linear map, and has exactly the same
values as $f$ at all the zeros. Likewise, for every
$(a_\rho)\in s(Z)$ the series $\sum_\rho a_\rho h_\rho$
converges in $\mathcal E$ and depends continuously on that
sequence. These assertions use (9) and choose the coefficient
decay power separately for each target seminorm.

\emph{Division on vertical strips.}
We now show the ideal equality, which does not follow merely
from constructing the cardinal functions.
Condition (1) implies a uniform, fixed-strip minimum-modulus
estimate
\[
 |L(\sigma+it,\eta)|
   \ge c_A(1+|t|)^{-C_A}
            \min\{1,\operatorname{dist}(\sigma+it,Z)\},
 \qquad |\sigma|\le A,\quad |t|\ge T_A.                 \tag{11}
\]
Here is a complete proof. Put $r=A+2$ and center a disk at
$c_t=2+it$. Choose $R\in[4r+4,4r+5]$ with zero-free boundary.
At large $|t|$ there are no poles or trivial zeros in that disk.
Its $M=O_A(\log(2+|t|))$ zeros are all simple and on the
critical line. Remove them by the finite Blaschke product
$B$ for this disk, as in MT-M11, and put $G=L/B$.
Polynomial strip growth gives $|G|\le C(1+|t|)^C$ on the
whole disk, while the Euler lower bound at $c_t$ and Harnack's
inequality give
$|G|\ge c(1+|t|)^{-C}$ on the disk of radius $R/2$.
The Harnack factor on that smaller disk is at most three.
The point $s=\sigma+it$ is within radius $r<R/4$.

If the disk has no zeros, this already proves (11).
Otherwise choose a zero $\rho_0$ in it nearest to $s$.
If $|\rho_0-c_t|>R/2$, every zero in the disk is at
distance at least $R/4$ from $s$. Each Blaschke factor
therefore has modulus at least $1/8$, and
$|B(s)|\ge8^{-M}$, a polynomial lower bound.
If $|\rho_0-c_t|\le R/2$, simplicity gives
$L'(\rho_0)=G(\rho_0)B'(\rho_0)$.
In normalized coordinates $|a_0|\le1/2$, the derivative
of its own Blaschke factor is at most $4/(3R)$, and
every other factor at $\rho_0$ is at most
$2|\rho_0-\rho_j|/R$. Thus
\[
 |B'(\rho_0)|\le C_A^M
                     \prod_{\rho_j\ne\rho_0}|\rho_0-\rho_j|.
\]
The assumed lower bound for $L'$, the upper bound for $G$,
and $M=O_A(\log(2+|t|))$ imply
\[
          \prod_{\rho_j\ne\rho_0}|\rho_0-\rho_j|
                          \ge c(1+|t|)^{-C}.           \tag{12}
\]
Since $\rho_0$ is nearest to $s$,
$|\rho_0-\rho_j|\le2|s-\rho_j|$. Each denominator in a
Blaschke factor at $s$ is at most two. Hence
\[
 |B(s)|\ge(4R)^{-M}|s-\rho_0|
                       \prod_{\rho_j\ne\rho_0}|\rho_0-\rho_j|,
\]
and (12) proves (11). Zeros outside the large disk have
distance greater than one from $s$, so using the distance
to the full set $Z$ makes no difference in the minimum.
The estimate at a zero is interpreted as zero on both
sides and follows by continuity.

\emph{The closed ideal is the full zero kernel.}
The inclusion
$\overline{\psi_\eta*\mathcal E}\subset K_Z$ is immediate
from multiplication of Mellin transforms.
Take $f\in K_Z$ and write $F=\mathcal Mf$.
The quotient $F/\Xi_\eta$ is entire because all zeros
are simple and have been cancelled. On every fixed strip,
(11) and the rapid vertical decay of $F$ and $F'$ imply
rapid decay of $F/L$ for large $|t|$. If the distance to
$Z$ is at least one, divide the rapid bound for $F$
by (11). If it is less than one, choose a nearest zero
$\rho$ and use
$F(s)=\int_\rho^s F'(w)\,dw$ on the straight segment;
that segment lies in a slightly larger fixed strip and
a bounded-height window. The distance factor cancels
the corresponding factor in (11).

Inverting the gamma factors now costs at most
$C_A(1+|t|)^{C_A}e^{\kappa|t|}$ on the strip.
For any $\epsilon>0$ the entire function
\[
             J_\epsilon(s)=e^{\epsilon s^2}
                                  \frac{F(s)}{\Xi_\eta(s)}
\]
therefore decays faster than every power on every vertical
strip, in fact with a Gaussian bound. On compact rectangles
it is bounded by entireness. Contour shifting and Fourier
inversion give $j_\epsilon\in\mathcal E$ with this Mellin
transform. Uniqueness of the transform gives
\[
                  \psi_\eta*j_\epsilon=g_\epsilon*f.
\]
Finally $g_\epsilon*f\to f$ in $\mathcal E$ as
$\epsilon\downarrow0$. To verify this in each $p_{N,j}$,
write the difference as the integral of
$g_\epsilon(y)(f^{(j)}(x-y)-f^{(j)}(x))$.
For $|y|\le1$, the weighted mean-value bound is
$|y|e^{N|y|}p_{N,j+1}(f)$; outside any fixed
neighborhood of zero, use
$(e^{N|y|}+1)p_{N,j}(f)$ and the Gaussian tail.
Both integrated bounds tend to zero.
Thus $f$ belongs to the closed convolution ideal,
proving (MT-M12a).

The coordinate map $Q_\eta\to s(Z)$ is now injective.
It is continuous by the integration-by-parts bounds above,
which descend to quotient seminorms because all coordinates
annihilate the ideal. The continuous series map
$(a_\rho)\mapsto[\sum_\rho a_\rho h_\rho]$ is its inverse.
This proves (MT-M12b), and the usual coordinate vectors
give (2), with representatives $h_\rho$.
Statement (2) implies (3). Statement (3) implies (1)
by MT-M11 with $m=1$, extending its high-zero estimate
over the finitely many remaining zeros. This completes
the equivalence.
\end{proof}

Replacing $\psi_\eta$ by
$g_b*\psi_\eta$ for any $b>0$ leaves the closed ideal unchanged:
Gaussian convolution has dense range by MT-M06 and convolution
by $\psi_\eta$ is continuous.

The theorem settles a precise conditional interpolation and
spectral-synthesis interface, rather than identifying every
quotient occurring in Masullo's draft. For this canonical
completed-Hecke quotient with simple zeros, the Schauder and
absolute-Schauder existence gates are exactly the polynomial
reciprocal-derivative gate. The extension to arbitrary
multiplicities requires bounds for the whole local Hermite
interpolation data and is not asserted here.

\subsection*{MT-M13. Nuclear arithmetic convolution under the exact basis gate}
Under the equivalent conditions of MT-M12, convolution by
$f\in\mathcal E$ acts on $Q_\eta\simeq s(Z)$ as
\[
            (a_\rho)\longmapsto
                         (\mathcal Mf(\rho)a_\rho).
\]
It has an explicit nuclear representation on that Fr\'echet
space and a trace-class extension to $\ell^2(Z)$, with
\[
                   \operatorname{Tr}C_f
                         =\sum_{\rho\in Z}\mathcal Mf(\rho).
                                                               \tag{MT-M13}
\]
Here $\operatorname{Tr}$ denotes the trace of that Hilbert-space extension.
The sum converges absolutely. For
$h^\sharp(x)=e^{-x}\overline{h(-x)}$ and
$f=h*h^\sharp$, this Hilbert-space operator is positive and
\[
                   \operatorname{Tr}C_f
                         =\sum_{\rho\in Z}|\mathcal Mh(\rho)|^2
                         \ge0.
\]

\begin{proof}
Convolution descends because it preserves the closed
convolution ideal. Multiplication of Mellin transforms
gives the diagonal action. The eigenvalues decrease
faster than every power of the ordinate, so their
absolute sum is finite by NT-Z10. A diagonal operator
with these eigenvalues on the orthonormal coordinate
basis of $\ell^2(Z)$ is trace class and has the stated
trace.

To see nuclearity without an unspecified basis estimate,
write $b_\rho=[h_\rho]$ and
$\delta_\rho([g])=\mathcal Mg(\rho)$. The functionals
$\delta_\rho$ are equicontinuous: for a fixed
$N>1/2$, their absolute values are uniformly bounded
by a constant times the quotient seminorm of $p_{N,0}$.
The representation
\[
 C_f=\sum_{\rho\in Z}(1+|\operatorname{Im}\rho|)^{-2}
       \delta_\rho\otimes
       \bigl((1+|\operatorname{Im}\rho|)^2
                                  \mathcal Mf(\rho)b_\rho\bigr)
\]
has summable scalar coefficients, an equicontinuous
family of functionals, and a bounded family of vectors
in every quotient seminorm by (9) and rapid decay.
It is therefore an explicit nuclear representation.
Its diagonal sum is the same absolutely convergent
trace displayed above; the unambiguous Hilbert
trace needs no additional abstract trace convention.

The weighted involution preserves $\mathcal E$ and satisfies
\[
      \mathcal M h^\sharp(s)
                =\overline{\mathcal Mh(1-\overline s)}.
\]
At every $\rho$ on the critical line,
$1-\overline\rho=\rho$. Thus the diagonal eigenvalues
for $h*h^\sharp$ are the nonnegative squares stated,
which proves positivity and its trace formula.
\end{proof}

This is a rigorous operator model conditional on the exact
simple-zero basis gate. Its Hilbert completion is constructed
from the zero coordinates; it is not an independently
constructed geometric or prime-side operator proving GRH.
Its location in the programme is Masullo 9.3.2.5--9.3.3,
the Tate-adjoint involution, and cote~1's nuclear trace
framework as carried in SA23 geo-H-lidskii/geo-H-genus.
